\FigureLayoutDeclare{img/1977/guangxi_liberal/q2_solution_fig1.png}{5.4cm}{0bp 0bp 62bp 0bp}

\chapter{1977年广西卷（文）}

\section{解答题}

\begin{problem}
计算：

\begin{enumerate}
\item \( 5\frac{3}{4}-2\left[ \left( \frac{1}{2} \right)^2-3\times\frac{3}{4} \right]\div\frac{2}{5} \).
\item 分解 \( x^3-5x^2+4x \) 的因式.
\item 化简 \( \frac{1}{\sin^2 \alpha}-\cot^2 \alpha \).
\item 计算 \( \cot 135^\circ+\cos \left( -60^\circ \right)-\tan^2 60^\circ+\sin 30^\circ \).
\item 不查表求 \(\lg\frac{3}{7}+\lg70-\lg3\) 的值.
\end{enumerate}
\end{problem}

\begin{solution}
\begin{enumerate}
\item 先计算括号内的值，并把除以 \(\frac{2}{5}\) 改写为乘以 \(\frac{5}{2}\)，得
\[
\begin{gathered}
5\frac{3}{4}-2\left[ \left( \frac{1}{2} \right)^2-3\times\frac{3}{4} \right]\div\frac{2}{5}
=5\frac{3}{4}-2\left(\frac{1}{4}-\frac{9}{4}\right)\times\frac{5}{2}\\
=5\frac{3}{4}-2\times(-2)\times\frac{5}{2}=5\frac{3}{4}+10=15\frac{3}{4}
\end{gathered}
\]

\item 先提取公因式 \(x\)，再分解所得二次三项式. 因为 \(1+4=5\)，且 \(1\times4=4\)，所以
\[
x^3-5x^2+4x=x\left(x^2-5x+4\right)=x\left(x-1\right)\left(x-4\right)
\]

\item 原式有意义时 \(\sin\alpha\ne0\)，且 \(\cot\alpha=\frac{\cos\alpha}{\sin\alpha}\). 利用 \(1-\cos^2\alpha=\sin^2\alpha\)，得
\[
\frac{1}{\sin^2\alpha}-\cot^2\alpha
=\frac{1}{\sin^2\alpha}-\frac{\cos^2\alpha}{\sin^2\alpha}
=\frac{1-\cos^2\alpha}{\sin^2\alpha}
=\frac{\sin^2\alpha}{\sin^2\alpha}=1
\]

\item 由三角函数的诱导公式和特殊角的三角函数值，\(\cot135^\circ=-1\)，\(\cos(-60^\circ)=\cos60^\circ=\frac{1}{2}\)，\(\tan60^\circ=\sqrt{3}\)，\(\sin30^\circ=\frac{1}{2}\). 因此
\[
\cot135^\circ+\cos(-60^\circ)-\tan^2 60^\circ+\sin30^\circ
=-1+\frac{1}{2}-\left(\sqrt{3}\right)^2+\frac{1}{2}=-3
\]

\item 由于各对数的真数均为正数，可以使用对数的和、差公式. 又因为 \(70=7\times10\)，所以
\[
\begin{gathered}
\lg\frac{3}{7}+\lg70-\lg3
=\left(\lg3-\lg7\right)+\left(\lg7+\lg10\right)-\lg3\\
=\lg10=1
\end{gathered}
\]
\end{enumerate}
\end{solution}

\begin{problem}
\(P\) 是圆外一点，过 \(P\) 作两条直线，一条交圆于 \(C\)，\(A\) 两点，一条交圆于 \(D\)，\(B\) 两点，求证 \( PA\cdot PC=PB\cdot PD \).
\end{problem}

\begin{solution}
\solutionfigure{\bitmapfigure[max width=4.2cm]{img/1977/guangxi_liberal/q2_solution_fig1.png}}

由题意，点 \(P\)，\(D\)，\(B\) 共线，点 \(P\)，\(C\)，\(A\) 共线. 因为 \(PA\) 与 \(PC\) 在同一直线上，\(PB\) 与 \(PD\) 在同一直线上，所以 \(\angle DPA=\angle CPB\). 又因为 \(P\)，\(C\)，\(A\) 共线且 \(C\) 在线段 \(PA\) 上，\(P\)，\(D\)，\(B\) 共线且 \(D\) 在线段 \(PB\) 上，所以 \(\angle PAD=\angle CAD\)，\(\angle PBC=\angle DBC\). \(\angle CAD\) 与 \(\angle DBC\) 都是同一圆中对着弦 \(CD\) 的圆周角，故 \(\angle PAD=\angle PBC\).

因此由两角分别相等，得 \(\triangle PDA\sim\triangle PCB\)，对应关系为 \(P\leftrightarrow P\)，\(A\leftrightarrow B\)，\(D\leftrightarrow C\). 于是
\[
\frac{PA}{PB}=\frac{PD}{PC}
\]
由于各线段长度均为正数，交叉相乘得
\[
PA\cdot PC=PB\cdot PD
\]
\end{solution}

\begin{problem}
\(m\) 是什么值的时候，方程 \(3x^2-\left(4m-1\right)x+m^2+\frac{3}{4}=0\) 有相等的两个实数根.
\end{problem}

\begin{solution}
这是关于 \(x\) 的一元二次方程，二次项系数为 \(3\ne0\). 它有相等的两个实数根，当且仅当判别式等于零. 方程的系数为 \(a=3\)，\(b=-(4m-1)=1-4m\)，\(c=m^2+\frac{3}{4}\)，所以
\[
\begin{gathered}
\Delta=b^2-4ac
=(1-4m)^2-12\left(m^2+\frac{3}{4}\right)\\
=16m^2-8m+1-12m^2-9\\
=4m^2-8m-8=4\left(m^2-2m-2\right)
\end{gathered}
\]
令 \(\Delta=0\)，得
\[
m^2-2m-2=0\quad\Longleftrightarrow\quad\left(m-1\right)^2=3
\]
因此 \(m=1+\sqrt{3}\) 或 \(m=1-\sqrt{3}\). 这两个值都使 \(\Delta=0\)，且二次项系数仍为 \(3\)，故方程确实有相等的两个实数根.
\end{solution}

\begin{problem}
波尔多液由胆矾，生石灰和水配制成，已知胆矾和生石灰的重量比是 \(1:2\)，生石灰和水的重量比是 \(1:120\)，问配制 \(486\text{ 斤}\) 波尔多液需胆矾，生石灰和水各多少斤？
\end{problem}

\begin{solution}
设胆矾的重量为 \(x\text{ 斤}\)，生石灰的重量为 \(y\text{ 斤}\)，水的重量为 \(z\text{ 斤}\). 由胆矾与生石灰的重量比为 \(1:2\)，得 \(y=2x\). 由生石灰与水的重量比为 \(1:120\)，得 \(z=120y=240x\).

三种物质的总重量为 \(486\text{ 斤}\)，所以 \(x+y+z=x+2x+240x=486\)，从而 \(243x=486\)，解得 \(x=2\). 于是 \(y=2x=4\)，\(z=240x=480\).

检验：\(2+4+480=486\)，且 \(2:4=1:2\)，\(4:480=1:120\)，符合题设条件.

答：需胆矾 \(2\text{ 斤}\)，生石灰 \(4\text{ 斤}\)，水 \(480\text{ 斤}\).
\end{solution}

\begin{problem}
已知桥的一孔为抛物线型（如图），跨度 \(AB=12\text{ 米}\)，桥孔高 \(PO=4\text{ 米}\)，建桥时每隔 \(2\text{ 米}\) 一个撑梁，求撑梁 \(C_1D_1\)，\(C_2D_2\)，\(C_3D_3\)，\(C_4D_4\) 各应取多长？

\bitmapfigure[max width=5.6cm]{img/1977/guangxi_liberal/q5_fig1.png}
\end{problem}

\begin{solution}
以抛物线的顶点 \(O\) 为坐标原点，以对称轴为 \(y\) 轴建立直角坐标系. 桥面所在的直线为 \(y=-4\)，且由于 \(AB=12\text{ 米}\)，所以 \(A\)、\(B\) 的坐标分别为 \((-6,-4)\)、\((6,-4)\).

设抛物线的标准方程为 \(x^2=-2py\)，其中 \(p>0\). 把点 \(B(6,-4)\) 代入，得 \(6^2=-2p\times(-4)\)，
所以 \(p=\frac{9}{2}\)，抛物线方程为
\[
x^2=-9y
\]
从而 \(y=-\frac{x^2}{9}\).

由图，四根撑梁的横坐标分别为 \(-4\)，\(-2\)，\(2\)，\(4\). 设撑梁底端在桥面 \(y=-4\) 上，横坐标为 \(t\) 的撑梁顶端在抛物线上，则顶端纵坐标为 \(-\frac{t^2}{9}\)，撑梁长度为
\[
\left(-\frac{t^2}{9}\right)-(-4)=4-\frac{t^2}{9}
\]

当 \(t=2\) 或 \(t=-2\) 时，撑梁长度为
\[
4-\frac{2^2}{9}=\frac{32}{9}=3\frac{5}{9}\text{ 米}
\]
当 \(t=4\) 或 \(t=-4\) 时，撑梁长度为
\[
4-\frac{4^2}{9}=\frac{20}{9}=2\frac{2}{9}\text{ 米}
\]

因此 \(C_1D_1=C_4D_4=2\frac{2}{9}\text{ 米}\)，\(C_2D_2=C_3D_3=3\frac{5}{9}\text{ 米}\).
\end{solution}
